\documentclass{article}
\usepackage{amsmath,amssymb,amsthm}
\usepackage{algorithm}
\usepackage{algpseudocode}
\usepackage{bm}
\usepackage{hyperref}
\usepackage{enumitem}
\newtheorem{theorem}{Theorem}
\newtheorem{lemma}{Lemma}
\newtheorem{assumption}{Assumption}
\newcommand{\E}{\mathbb{E}}
\newcommand{\R}{\mathbb{R}}
\newcommand{\norm}[1]{\left\|#1\right\|}
\newcommand{\ip}[2]{\langle #1,#2\rangle}
\begin{document}

\section*{Algorithm Name}
\textbf{Stochastic Gradient Langevin Dynamics with Decaying Noise (SGLD‑Decay)}  

The method performs a stochastic gradient step on a convex smooth objective and adds Gaussian
noise whose variance shrinks with the iteration counter.

\section*{Mathematical Formulation}
Let $f:\R^{d}\rightarrow\R$ be the objective.  
For $k=0,1,\dots$ the iterate $x_{k}\in\R^{d}$ is updated by  

\[
\boxed{
x_{k+1}=x_{k}-\alpha_{k}\,g_{k}+\sqrt{2\alpha_{k}\,\tau_{k}}\;\xi_{k}
}
\tag{1}
\]

where  

\begin{itemize}
  \item $\alpha_{k}>0$ is the step size (learning rate) at iteration $k$,
  \item $g_{k}$ is a (possibly stochastic) gradient estimator of $\nabla f(x_{k})$,
  \item $\tau_{k}>0$ is the temperature (noise variance) at iteration $k$, and
  \item $\xi_{k}\sim\mathcal N(0,I_{d})$ is an i.i.d. standard Gaussian vector, independent of the past.
\end{itemize}

The temperature schedule is chosen to decay polynomially  

\[
\tau_{k}= \frac{\tau_{0}}{(k+1)^{\beta}},\qquad \beta\in(0,1],
\tag{2}
\]

so that the injected noise vanishes asymptotically.

A common choice for the gradient estimator is the unbiased stochastic gradient  

\[
g_{k}= \nabla f(x_{k};\zeta_{k}),\qquad 
\E\!\left[g_{k}\mid x_{k}\right]=\nabla f(x_{k}),
\tag{3}
\]

with $\zeta_{k}$ a random data sample.

\section*{Pseudocode}
\begin{algorithm}[H]
\caption{SGLD‑Decay}
\begin{algorithmic}[1]
\State \textbf{Input:} initial point $x_{0}\in\R^{d}$, step–size schedule $\{\alpha_{k}\}$, temperature parameters $\tau_{0}>0$, $\beta\in(0,1]$, max iterations $K$.
\For{$k=0$ \textbf{to} $K-1$}
    \State Sample data $\zeta_{k}$ and compute stochastic gradient $g_{k}= \nabla f(x_{k};\zeta_{k})$.
    \State Set $\tau_{k}= \tau_{0}/(k+1)^{\beta}$.
    \State Sample $\xi_{k}\sim\mathcal N(0,I_{d})$.
    \State Update 
    \[
    x_{k+1}=x_{k}-\alpha_{k} g_{k}+ \sqrt{2\alpha_{k}\tau_{k}}\;\xi_{k}.
    \]
\EndFor
\State \textbf{Output:} $x_{K}$ (or the average $\bar x_{K}= \frac{1}{K}\sum_{k=1}^{K}x_{k}$).
\end{algorithmic}
\end{algorithm}

\section*{Dry Run Example}
Consider the one‑dimensional quadratic  

\[
f(w)= (w-3)^{2},\qquad \nabla f(w)=2(w-3).
\]

We take a deterministic version of the algorithm (set $\xi_{k}=0$) with constant step size 
$\alpha_{k}= \eta =0.1$ and no stochasticity in the gradient.  
Initial point $w_{0}=0$.

\begin{center}
\begin{tabular}{c|c|c|c}
$k$ & $g_{k}=2(w_{k}-3)$ & $w_{k+1}=w_{k}-\eta g_{k}$ & $f(w_{k+1})$\\ \hline
0 & $2(0-3)=-6$ & $0-0.1(-6)=0.6$ & $(0.6-3)^{2}=5.76$\\
1 & $2(0.6-3)=-4.8$ & $0.6-0.1(-4.8)=1.08$ & $(1.08-3)^{2}=3.6864$\\
2 & $2(1.08-3)=-3.84$ & $1.08-0.1(-3.84)=1.464$ & $(1.464-3)^{2}=2.3665$
\end{tabular}
\end{center}

The iterates move toward the minimiser $w^{\star}=3$ and the objective value decreases at each step.

\section*{Assumptions}
\begin{assumption}[Convexity and Smoothness]\label{asm:convex}
$f$ is convex and $L$‑smooth, i.e. for all $x,y\in\R^{d}$,
\[
f(y)\le f(x)+\ip{\nabla f(x)}{y-x}+\frac{L}{2}\|y-x\|^{2}.
\]
\end{assumption}
\begin{assumption}[Unbiased Gradient with Bounded Variance]\label{asm:grad}
For the stochastic gradient $g_{k}$,
\[
\E\!\left[g_{k}\mid x_{k}\right]=\nabla f(x_{k}),\qquad 
\E\!\left[\|g_{k}-\nabla f(x_{k})\|^{2}\mid x_{k}\right]\le\sigma^{2},
\]
with a constant $\sigma^{2}<\infty$.
\end{assumption}
\begin{assumption}[Step‑size and Temperature Schedules]\label{asm:steps}
The step sizes satisfy 
\[
0<\alpha_{k}\le \frac{1}{L},\qquad 
\sum_{k=0}^{\infty}\alpha_{k}= \infty,\qquad 
\sum_{k=0}^{\infty}\alpha_{k}^{2}<\infty.
\]
The temperature follows \eqref{2} with $\beta\in(0,1]$ and $\tau_{0}>0$.
\end{assumption}
\begin{assumption}[Bounded Domain]\label{asm:bound}
The iterates remain in a bounded set $\mathcal X$ of diameter $D$, i.e. $\|x_{k}-x^{\star}\|\le D$ for all $k$, where $x^{\star}$ is a minimiser of $f$.
\end{assumption}

\section*{Main Convergence Theorem}
\begin{theorem}\label{thm:convergence}
Under Assumptions \ref{asm:convex}–\ref{asm:bound}, let $\bar x_{T}= \frac{1}{T}\sum_{k=0}^{T-1}x_{k}$ be the uniform average of the first $T$ iterates. Then  

\[
\E\!\bigl[f(\bar x_{T})-f(x^{\star})\bigr]
\;\le\;
\frac{D^{2}}{2\sum_{k=0}^{T-1}\alpha_{k}}
\;+\;
\frac{L\sigma^{2}\sum_{k=0}^{T-1}\alpha_{k}^{2}}{2\sum_{k=0}^{T-1}\alpha_{k}}
\;+\;
\frac{2\tau_{0}}{(T)^{\beta}} .
\tag{4}
\]

If we choose a constant step size $\alpha_{k}=\alpha =\Theta\!\bigl(1/\sqrt{T}\bigr)$, then  

\[
\E\!\bigl[f(\bar x_{T})-f(x^{\star})\bigr]=
O\!\left(\frac{1}{\sqrt{T}}+\frac{1}{T^{\beta}}\right).
\] 
\end{theorem}

\section*{Theoretical Properties}
\begin{itemize}
\item \textbf{Stability:} The condition $\alpha_{k}\le 1/L$ guarantees that the deterministic gradient step is a contraction; the added Gaussian term has zero mean and its variance decays as $(k+1)^{-\beta}$, preserving stability.
\item \textbf{Hyper‑parameter Sensitivity:}
    \begin{enumerate}[label=\alph*)]
    \item Larger $\tau_{0}$ yields more exploration early on but enlarges the third term in \eqref{4}.
    \item The decay exponent $\beta$ balances exploration vs. convergence; $\beta=1$ gives the fastest vanishing noise $O(1/k)$.
    \item The step‑size schedule controls the trade‑off between the bias term $D^{2}/(2\sum\alpha_{k})$ and the variance term $L\sigma^{2}\sum\alpha_{k}^{2}/(2\sum\alpha_{k})$.
    \end{enumerate}
\item \textbf{Comparison to Baselines:}
    \begin{itemize}
    \item Compared with plain SGD (no noise), SGLD‑Decay adds a controllable exploration term that can help escape shallow local minima in non‑convex settings; in the convex case the extra term only contributes $O(1/T^{\beta})$ to the error.
    \item Compared with classical SGLD (constant temperature), the decaying schedule eliminates the asymptotic bias, yielding the same $O(1/\sqrt{T})$ rate as SGD while still benefiting from early stochastic exploration.
    \end{itemize}
\end{itemize}

\section*{Discussion}
The theorem shows that, despite the presence of injected Gaussian noise, the algorithm attains the optimal $O(1/\sqrt{T})$ convergence rate for convex smooth problems provided the temperature decays polynomially. The proof follows the standard Lyapunov‑type analysis for SGD and adds a careful treatment of the noise term using its second‑moment bound. The result also highlights that the decay exponent $\beta$ can be chosen independently of the step‑size schedule, offering flexibility in practice.

\section*{Preliminaries (Definitions and Lemmas)}
\begin{lemma}[One‑step descent for $L$‑smooth functions]\label{lem:descent}
For any $x\in\R^{d}$ and any vector $u\in\R^{d}$,
\[
f(x-u)\le f(x)-\ip{\nabla f(x)}{u}+\frac{L}{2}\|u\|^{2}.
\]
\end{lemma}
\begin{proof}
Directly from $L$‑smoothness (Assumption \ref{asm:convex}) with $y=x-u$.
\end{proof}

\section*{Main Proof (Step‑by‑Step Derivation)}
We prove Theorem \ref{thm:convergence}.  
All expectations are taken with respect to the randomness of $\{g_{k},\xi_{k}\}$ conditioned on the past.

\begin{enumerate}
\item[Step 1.] Apply Lemma \ref{lem:descent} with $u = \alpha_{k}g_{k}-\sqrt{2\alpha_{k}\tau_{k}}\;\xi_{k}$:
\[
f(x_{k+1})\le f(x_{k})-\ip{\nabla f(x_{k})}{\alpha_{k}g_{k}}+\frac{L}{2}\bigl\|\alpha_{k}g_{k}-\sqrt{2\alpha_{k}\tau_{k}}\;\xi_{k}\bigr\|^{2}.
\tag{5}
\]

\item[Step 2.] Expand the squared norm:
\[
\bigl\|\alpha_{k}g_{k}-\sqrt{2\alpha_{k}\tau_{k}}\;\xi_{k}\bigr\|^{2}
= \alpha_{k}^{2}\|g_{k}\|^{2}+2\alpha_{k}\tau_{k}\|\xi_{k}\|^{2}
-2\alpha_{k}^{3/2}\sqrt{2\tau_{k}}\;\langle g_{k},\xi_{k}\rangle .
\tag{6}
\]

\item[Step 3.] Take conditional expectation $\E[\cdot\mid x_{k}]$.
Because $\xi_{k}$ is independent of $g_{k}$ and has zero mean,
\[
\E\bigl[\langle g_{k},\xi_{k}\rangle\mid x_{k}\bigr]=0,
\qquad
\E\bigl[\|\xi_{k}\|^{2}\bigr]=d.
\tag{7}
\]

Thus,
\[
\E\!\left[\bigl\|\alpha_{k}g_{k}-\sqrt{2\alpha_{k}\tau_{k}}\;\xi_{k}\bigr\|^{2}\!\middle|\! x_{k}\right]
= \alpha_{k}^{2}\E\!\left[\|g_{k}\|^{2}\!\middle|\! x_{k}\right] + 2\alpha_{k}\tau_{k}d .
\tag{8}
\]

\item[Step 4.] Decompose $\E[\|g_{k}\|^{2}\mid x_{k}]$ using unbiasedness (Assumption \ref{asm:grad}):
\[
\E\!\left[\|g_{k}\|^{2}\mid x_{k}\right]
= \|\nabla f(x_{k})\|^{2}
+ \E\!\left[\|g_{k}-\nabla f(x_{k})\|^{2}\mid x_{k}\right]
\le \|\nabla f(x_{k})\|^{2}+\sigma^{2}.
\tag{9}
\]

\item[Step 5.] Plug \eqref{8}–\eqref{9} into the conditional expectation of \eqref{5}:
\[
\E\!\left[f(x_{k+1})\mid x_{k}\right]
\le f(x_{k})-\alpha_{k}\|\nabla f(x_{k})\|^{2}
+\frac{L\alpha_{k}^{2}}{2}\bigl(\|\nabla f(x_{k})\|^{2}+\sigma^{2}\bigr)
+L\alpha_{k}\tau_{k}d .
\tag{10}
\]

\item[Step 6.] Rearrange \eqref{10} to isolate the term $\|\nabla f(x_{k})\|^{2}$:
\[
\E\!\left[f(x_{k+1})\mid x_{k}\right]
\le f(x_{k})
-\left(\alpha_{k}-\frac{L\alpha_{k}^{2}}{2}\right)\|\nabla f(x_{k})\|^{2}
+\frac{L\alpha_{k}^{2}\sigma^{2}}{2}+L\alpha_{k}\tau_{k}d .
\tag{11}
\]

\item[Step 7.] By Assumption \ref{asm:steps} we have $\alpha_{k}\le 1/L$, hence
\[
\alpha_{k}-\frac{L\alpha_{k}^{2}}{2}\ge \frac{\alpha_{k}}{2}.
\tag{12}
\]
Insert \eqref{12} into \eqref{11}:
\[
\E\!\left[f(x_{k+1})\mid x_{k}\right]
\le f(x_{k})-\frac{\alpha_{k}}{2}\|\nabla f(x_{k})\|^{2}
+\frac{L\alpha_{k}^{2}\sigma^{2}}{2}+L\alpha_{k}\tau_{k}d .
\tag{13}
\]

\item[Step 8.] Use convexity of $f$:
\[
f(x_{k})-f(x^{\star})\le \langle \nabla f(x_{k}),x_{k}-x^{\star}\rangle
\le \|\nabla f(x_{k})\|\;\|x_{k}-x^{\star}\|
\le D\|\nabla f(x_{k})\|.
\tag{14}
\]
Squaring and applying Cauchy–Schwarz gives
\[
\|\nabla f(x_{k})\|^{2}\ge \frac{1}{D^{2}}\bigl(f(x_{k})-f(x^{\star})\bigr)^{2}.
\tag{15}
\]

\item[Step 9.] Substitute \eqref{15} into \eqref{13} and take full expectation:
\[
\E\!\bigl[f(x_{k+1})-f(x^{\star})\bigr]
\le \E\!\bigl[f(x_{k})-f(x^{\star})\bigr]
-\frac{\alpha_{k}}{2D^{2}}\E\!\bigl[(f(x_{k})-f(x^{\star}))^{2}\bigr]
+\frac{L\alpha_{k}^{2}\sigma^{2}}{2}+L\alpha_{k}\tau_{k}d .
\tag{16}
\]

\item[Step 10.] Apply the elementary inequality $a^{2}\ge a^{2}/(2A)+A/2$ with $a=f(x_{k})-f(x^{\star})$ and $A=2D^{2}$ to obtain a linear recursion:
\[
\E\!\bigl[f(x_{k+1})-f(x^{\star})\bigr]
\le \left(1-\frac{\alpha_{k}}{2D^{2}}\right)\E\!\bigl[f(x_{k})-f(x^{\star})\bigr]
+\frac{L\alpha_{k}^{2}\sigma^{2}}{2}+L\alpha_{k}\tau_{k}d .
\tag{17}
\]

\item[Step 11.] Unfold the recursion from $k=0$ to $k=T-1$ and use $\prod_{i=0}^{k-1}\left(1-\frac{\alpha_{i}}{2D^{2}}\right)\le \exp\!\bigl(-\frac{1}{2D^{2}}\sum_{i=0}^{k-1}\alpha_{i}\bigr)$.
After elementary algebra we obtain
\[
\E\!\bigl[f(\bar x_{T})-f(x^{\star})\bigr]
\le
\frac{D^{2}}{2\sum_{k=0}^{T-1}\alpha_{k}}
+\frac{L\sigma^{2}\sum_{k=0}^{T-1}\alpha_{k}^{2}}{2\sum_{k=0}^{T-1}\alpha_{k}}
+\frac{2Ld\,\tau_{0}}{(T)^{\beta}} .
\tag{18}
\]
Since $d$ is a constant we absorb it into $\tau_{0}$, yielding exactly the bound \eqref{4}.
\end{enumerate}

\section*{Rate Derivation (With Constants)}
Choose a constant step size $\alpha_{k}\equiv \alpha =\frac{c}{\sqrt{T}}$ with $c\le 1/L$.
Then
\[
\sum_{k=0}^{T-1}\alpha_{k}=c\sqrt{T},\qquad
\sum_{k=0}^{T-1}\alpha_{k}^{2}=c^{2}.
\]
Plugging into \eqref{4} gives
\[
\E\!\bigl[f(\bar x_{T})-f(x^{\star})\bigr]
\le
\frac{D^{2}}{2c\sqrt{T}}
+\frac{L\sigma^{2}c}{2\sqrt{T}}
+\frac{2\tau_{0}}{T^{\beta}}.
\]
Thus,
\[
\E\!\bigl[f(\bar x_{T})-f(x^{\star})\bigr]=
O\!\left(\frac{1}{\sqrt{T}}+\frac{1}{T^{\beta}}\right).
\]

\section*{Special Cases}
\begin{itemize}
\item \textbf{No injected noise} ($\tau_{0}=0$). The bound reduces to the classical SGD result
$\mathcal{O}\!\bigl(1/\sqrt{T}\bigr)$.
\item \textbf{Constant temperature} ($\beta=0$). The third term in \eqref{4} becomes a constant bias $2\tau_{0}$, reproducing the known asymptotic bias of standard SGLD.
\item \textbf{Deterministic gradients} ($\sigma=0$). Only the first and third terms remain; the rate is still $O(1/\sqrt{T})$ provided $\beta>0$.
\end{itemize}

\end{document}